{\bf   Computing subfields in pure trascendental extensions}

\medskip

\centerline{Jaime Gutierrez and Rosario Rubio}

\centerline{Universidad de Cantabria, Spain}

\bigskip


\noindent  Rational function fields
$K(x_1,\cdots,x_n)$ arise in various contexts
within mathematics
and computer science. Two  examples are
the invariant theory  and the design of diffractive
optical systems. Several questions related to
subfields of $K(x_1,\cdots,x_n)$ occur in control
theory in the structural
identifiability problem, i.e., to check whether a differential equation model
process uniquely determines the internal parameters of the process
actions. Multivariate decompostion of
rational functions proves useful for inverting rational
functions. \par

\noindent The motivation of this paper on the one hand is
to generalize the univariate rational function decomposition
problem to the multivariate one; on the other hand, we want to
give new algorithms for questions concerning unirational and
rational fields. We also hope that our approach will be a step to
the better understanding of the design of optical systems for
computation. For this we consider the
following computational problem:

\medskip

\noindent Given $f_1,\cdots,f_m\in
K(x_1,\cdots,x_n)$, we want to know if there
exist
$h_1,\cdots,h_t\in K(x_1,\cdots,x_n)$ such that
$K(f_1,\cdots,f_m)\subset K(h_1,\cdots,h_t)$,
is a proper  algebraic extension.

\noindent And in the affirmative case,
compute such
$h$'s and rational functions
$g_1,\cdots,g_m\in K(y_1,\cdots,y_t)$ such
that

$$f_i=g_i(h_1,\cdots,h_t) {\rm \hskip 2mm for
\hskip 3mm all\hskip 2mm} i\in\{1,\cdots,m\}$$



\medskip

\noindent In this talk we present a polynomial time
method to solve the above problem.
The algorithm is based on
the computation of subfields (using blocks of decomposition) of
an algebraic extension. The basic idea underlying our approach
is the translation of decompositions issues to factorization
over algebraic extensions and then we find minimal blocks.


\end
